\documentclass[11pt]{article}

\usepackage{amsmath,amssymb,amsthm,mathtools}
\usepackage[margin=1in]{geometry}
\usepackage{enumitem}

% Theorem environments
\newtheorem{theorem}{Theorem}[section]
\newtheorem{lemma}[theorem]{Lemma}
\newtheorem{corollary}[theorem]{Corollary}
\newtheorem{definition}[theorem]{Definition}
\newtheorem{remark}[theorem]{Remark}

% Notation shortcuts
\newcommand{\R}{\mathbb{R}}
\newcommand{\N}{\mathbb{N}}
\newcommand{\Rp}{\R_+}
\newcommand{\Simplenat}{\mathrm{Simple}_{\mathrm{nat}}}
\newcommand{\Simpleord}{\mathrm{Simple}_{\preceq}}

\title{Higher-Dimensional Bootstrap for the Simplicity Conjecture:\\
Finite Product Menus}
\author{Draft --- not for circulation}
\date{\today}

\begin{document}

\maketitle

\begin{abstract}
We prove that on a finite product menu $Y = \prod_{i=1}^k M_i \subseteq [0,1]^k$,
a canonical pricing function $p_0$ (convex, coordinatewise nondecreasing, $p_0(0)=0$)
satisfies the simplicity property under the natural numerical order if and only if
$p_0$ is additively separable on $Y$.  We further prove an existential
order-normalization result: if simplicity holds for \emph{some} family of total
orders on the coordinate sets, then it also holds under the natural order,
and hence the price is additively separable.
The argument reduces the $k$-good problem to the two-good
characterization (Theorem~\ref{thm:k2}) via a slice-reduction lemma that uses
pointwise supporting-slope pinning of frozen coordinates.
\end{abstract}

%% ===================================================================
\section{Model and Definitions}\label{sec:model}
%% ===================================================================

\paragraph{Environment.}
One seller offers $k \geq 1$ goods to a single buyer.
The buyer has an additive valuation $\theta = (\theta_1,\dots,\theta_k) \in \Rp^k$
and evaluates a bundle (allocation vector) $y \in [0,1]^k$ at $\theta \cdot y = \sum_{i=1}^k \theta_i y_i$.

\paragraph{Mechanisms and canonical pricing.}
A direct mechanism $(x,t)$ consists of an allocation rule $x\colon \Rp^k \to [0,1]^k$
and a payment rule $t\colon \Rp^k \to \R$.
We impose incentive compatibility (IC), individual rationality (IR), and
no positive transfers (NPT): $t(\theta) \geq 0$ for all $\theta \in \Rp^k$.
Under these conditions, the buyer's indirect utility $U(\theta) = \theta \cdot x(\theta) - t(\theta)$
is convex and nondecreasing, and there exists a \emph{canonical pricing function}
\[
  p_0\colon [0,1]^k \to [0,\infty],
\]
which is the Fenchel conjugate of $U$, satisfying:
\begin{enumerate}[label=(\roman*)]
  \item $p_0$ is convex and lower semicontinuous (closed),
  \item $p_0$ is coordinatewise nondecreasing,
  \item $p_0(0) = 0$,
  \item $U(\theta) = \max_{y \in [0,1]^k}\bigl[\theta \cdot y - p_0(y)\bigr]$.
\end{enumerate}

\paragraph{Finite product menu.}
For each coordinate $i \in [k] \coloneqq \{1,\dots,k\}$, let
\[
  M_i = \{m_{i,1} < m_{i,2} < \cdots < m_{i,n_i}\} \subseteq [0,1], \qquad n_i \geq 2,
\]
and define the finite product menu
\[
  Y \coloneqq \prod_{i=1}^k M_i.
\]
We assume $p_0(y) < \infty$ for every $y \in Y$.

\paragraph{Utility on $Y$.}
For each type $\theta \in \Rp^k$ and bundle $y \in Y$,
\[
  u_\theta(y) \coloneqq \theta \cdot y - p_0(y) = \sum_{i=1}^k \theta_i y_i - p_0(y).
\]

\paragraph{Immediate-adjacency neighbors.}
Given a family of total orders $\preceq = (\preceq_i)_{i=1}^k$ on the sets $M_i$,
define the set of \emph{immediate $\preceq$-neighbors} of $y \in Y$ by
\[
  N_{\preceq}(y) \coloneqq \bigl\{z \in Y : \exists\, i \in [k] \;\text{such that}\;
  z_{-i} = y_{-i} \;\text{and}\;
  z_i \;\text{is the immediate predecessor or successor of}\; y_i \;\text{in}\; \preceq_i\bigr\}.
\]
When $\preceq_i$ is the natural numerical order $\leq$ on $M_i$ for every $i$,
we write $N_{\mathrm{nat}}(y)$.

\begin{definition}[Simplicity]\label{def:simplicity}
The ordered menu $(Y, p_0|_Y, \preceq)$ is \emph{simple} if
\[
  \forall\, \theta \in \Rp^k\;\; \forall\, y \in Y, \quad
  u_\theta(y) < \max_{w \in Y} u_\theta(w) \;\;\Longrightarrow\;\;
  \exists\, z \in N_{\preceq}(y) \;\text{with}\; u_\theta(z) > u_\theta(y).
\]
We write $\Simpleord(Y, p_0)$ for this property and $\Simplenat(Y, p_0)$
when the natural numerical orders are used.
\end{definition}

\begin{definition}[Broad simplicity]\label{def:broad-simplicity}
Say that $(Y, p_0|_Y)$ is \emph{broadly simple} if, for every
$\theta \in \Rp^k$ and every $y \in Y$ that is not a global maximizer
of $u_\theta$, there exist a coordinate $i \in [k]$ and a value
$t \in M_i$ such that $u_\theta(y_{-i}, t) > u_\theta(y)$.
Equivalently, every non-maximizer can be improved by changing a single
coordinate to \emph{some} value in that coordinate's menu (not
necessarily an adjacent one).
\end{definition}

\begin{lemma}[Equivalence of neighbor notions under convexity]%
\label{lem:broad-narrow}
When $p_0$ is convex, $\Simplenat(Y, p_0)$ and broad simplicity
of $(Y, p_0|_Y)$ are equivalent.
\end{lemma}

\begin{proof}
$\Simplenat$ trivially implies broad simplicity, since every
natural-adjacent neighbor is in particular a one-coordinate change.

Conversely, suppose $(Y, p_0|_Y)$ is broadly simple and fix
$\theta \in \Rp^k$, a non-maximizer $y$, and a coordinate~$i$ and
value $t \in M_i$ with $u_\theta(y_{-i}, t) > u_\theta(y)$.
Define $\phi(s) \coloneqq \theta_i\, s - p_0(y|_{i \to s})$ for
$s \in M_i$.  Since $p_0$ is convex, $\phi$ is discretely concave
on $M_i$ (Lemma~\ref{lem:concavity}).  We have
$\phi(t) > \phi(y_i)$, so $y_i$ is not a global maximizer of~$\phi$.
By Lemma~\ref{lem:concavity}, every local maximizer of a discretely
concave function is global, so $y_i$ is not a local maximizer either.
Hence some natural-adjacent $s' \in M_i$ satisfies
$\phi(s') > \phi(y_i)$, which gives $u_\theta(y_{-i}, s') > u_\theta(y)$.
Thus $y$ has an improving natural-adjacent neighbor, and $\Simplenat$
holds.
\end{proof}

\begin{remark}\label{rem:model-pdf}
The definition of ``neighbor'' in the original model
allows any one-coordinate change in a total order, which on a finite
set coincides with our broad-simplicity notion.
Lemma~\ref{lem:broad-narrow} therefore shows that, in the convex
canonical-pricing setting, our immediate-adjacency theorems
(Theorems~\ref{thm:k2}--\ref{thm:order-norm}) also hold under the
original model's neighbor definition.
\end{remark}

\paragraph{Two-coordinate slices and slice price matrices.}
For distinct $i,j \in [k]$ and an anchor $a \in \prod_{\ell \neq i,j} M_\ell$,
define the \emph{two-coordinate slice}
\[
  S_{ij}(a) \coloneqq \bigl\{y \in Y : y_{-\{i,j\}} = a \bigr\}.
\]
The associated \emph{slice price matrix} is
\[
  P^{ij,a}_{r,s} \coloneqq p_0(y) \quad \text{where } y_i = m_{i,r},\; y_j = m_{j,s},\; y_{-\{i,j\}} = a,
\]
for $r \in [n_i]$, $s \in [n_j]$.

%% ===================================================================
\section{Statement of Results}\label{sec:results}
%% ===================================================================

\begin{theorem}[Two-good characterization]\label{thm:k2}
Let $Y = M_1 \times M_2$ where $M_1, M_2 \subseteq [0,1]$ are finite sets
with their natural orders, and let $p_0|_Y$ be the restriction to $Y$ of a
convex nondecreasing function on $[0,1]^2$ with $p_0(0,0) = 0$.  Then
\[
  \Simplenat(Y, p_0) \quad\Longleftrightarrow\quad
  p_0 \;\text{is additively separable on}\; Y,
\]
i.e., there exist functions $p_1\colon M_1 \to \R$ and $p_2\colon M_2 \to \R$
such that $p_0(y_1,y_2) = p_1(y_1) + p_2(y_2)$ for all $(y_1,y_2) \in Y$.
\end{theorem}

\begin{theorem}[Main theorem]\label{thm:main}
Let $k \geq 2$ and $Y = \prod_{i=1}^k M_i \subseteq [0,1]^k$ with $p_0$ the
canonical pricing function of an IC/IR/NPT mechanism, finite on $Y$.  Then
\[
  \Simplenat(Y, p_0) \quad\Longleftrightarrow\quad
  \exists\, p_1,\dots,p_k \;(p_i\colon M_i \to \R) \;\;\text{such that}\;\;
  p_0(y) = \sum_{i=1}^k p_i(y_i) \quad \forall\, y \in Y.
\]
\end{theorem}

\begin{theorem}[Order-Normalization Lemma]\label{thm:order-norm}
Under the same hypotheses,
\[
  \bigl(\exists\, \preceq = (\preceq_i)_{i=1}^k\;\; \Simpleord(Y, p_0)\bigr)
  \;\;\Longleftrightarrow\;\;
  \Simplenat(Y, p_0)
  \;\;\Longleftrightarrow\;\;
  p_0 \;\text{additively separable on}\; Y.
\]
\end{theorem}

%% ===================================================================
%% ===================================================================
\section{Proof of Theorem~\ref{thm:k2} (Two-Good Characterization)}\label{sec:k2proof}
%%  ===================================================================
%% -------------------------------------------------------------------

We prove Theorem~\ref{thm:k2} for $k=2$; write $M_1 = \{a_1 < \cdots < a_m\}$
and $M_2 = \{b_1 < \cdots < b_n\}$, and set $p_{rs} \coloneqq p_0(a_r, b_s)$.
For adjacent pairs $a_r < a_{r+1}$ and $b_s < b_{s+1}$, define the
\emph{elementary $2\times 2$ face} $F = \{(a_r,b_s),(a_{r+1},b_s),(a_r,b_{s+1}),(a_{r+1},b_{s+1})\}$
and the \emph{mixed difference}
\[
  \Delta_F \coloneqq p_{r+1,s+1} - p_{r+1,s} - p_{r,s+1} + p_{r,s}.
\]

\begin{lemma}[Slope interval pins one coordinate]\label{lem:slopepin}
Fix a point $y = (a_r, b_s) \in Y$.  For each coordinate $i \in \{1,2\}$, define
the left and right secant slopes of the one-dimensional section
$t \mapsto p_0(y|_{i \to t})$ at $y_i$ (when they exist).
Write $L_i(y)$ for the maximum left secant slope and $R_i(y)$ for the minimum
right secant slope, with the convention $L_i(y) = 0$ at a left boundary and
$R_i(y) = +\infty$ at a right boundary.  Then:
\begin{enumerate}[label=(\alph*)]
  \item $0 \leq L_i(y) \leq R_i(y)$, by convexity and monotonicity of $p_0$.
  \item Any $\alpha \in [L_i(y), R_i(y)]$ makes $y_i$ a maximizer of
    $t \mapsto \alpha\, t - p_0(y|_{i \to t})$ among all immediate neighbors
    of $y_i$ in $M_i$.
\end{enumerate}
\end{lemma}

\begin{proof}
The one-dimensional section is convex (nondecreasing secant slopes) and
nondecreasing, since $p_0$ is.  For any left neighbor $s < y_i$ and right
neighbor $t > y_i$, convexity gives
$\frac{p_0(y|_{i\to y_i}) - p_0(y|_{i\to s})}{y_i - s}
\leq \frac{p_0(y|_{i\to t}) - p_0(y|_{i\to y_i})}{t - y_i}$,
so $L_i(y) \leq R_i(y)$.
If $\alpha \in [L_i(y), R_i(y)]$, then $\alpha \geq L_i(y)$ makes every
leftward move non-improving:
$\alpha(y_i - s) \geq p_0(y|_{i\to y_i}) - p_0(y|_{i\to s})$;
and $\alpha \leq R_i(y)$ makes every rightward move non-improving:
$\alpha(t - y_i) \leq p_0(y|_{i\to t}) - p_0(y|_{i\to y_i})$.
\end{proof}

\begin{lemma}[Nonzero face creates a simplicity trap]\label{lem:trap}
If some elementary face $F$ has $\Delta_F \neq 0$, then there exist
$\theta \in \Rp^2$ and a face vertex $y$ such that $y$ is not a global
maximizer of $u_\theta$ on $Y$, yet every adjacent neighbor $z \in N_{\mathrm{nat}}(y)$
satisfies $u_\theta(z) \leq u_\theta(y)$.
\end{lemma}

\begin{proof}
Let $F$ involve the adjacent pair $a_r < a_{r+1}$ in coordinate~1 and
$b_s < b_{s+1}$ in coordinate~2.  Write the four vertices as
$y_{00} = (a_r, b_s)$, $y_{10} = (a_{r+1}, b_s)$,
$y_{01} = (a_r, b_{s+1})$, $y_{11} = (a_{r+1}, b_{s+1})$.

\smallskip
\emph{Case $\Delta_F > 0$.}
Set $y = y_{10} = (a_{r+1}, b_s)$ and choose
\[
  \theta_1 = \frac{p_{r+1,s} - p_{r,s}}{a_{r+1} - a_r} = L_1(y),
  \qquad
  \theta_2 = \frac{p_{r+1,s+1} - p_{r+1,s}}{b_{s+1} - b_s} = R_2(y).
\]
Both are nonneg\-ative by mono\-tonicity of $p_0$.
Note that $\theta_1 = L_1(y)$ is the left secant slope
in coordinate~1 at $y_{10}$, and $\theta_2 = R_2(y)$ is the right secant slope
in coordinate~2 at $y_{10}$.
By Lemma~\ref{lem:slopepin}(b), $\theta_1 \in [L_1(y), R_1(y)]$ (taking
the left endpoint) pins coordinate~1: no adjacent move in coordinate~1
improves $u_\theta$.  Similarly $\theta_2 \in [L_2(y), R_2(y)]$ (right endpoint)
pins coordinate~2.  Therefore every neighbor of $y$ in $Y$ is non-improving.

Now check that $y$ is not a global maximizer.
At the diagonal vertex $y_{01} = (a_r, b_{s+1})$:
\begin{align*}
  u_\theta(y_{01}) - u_\theta(y_{10})
  &= \theta_1(a_r - a_{r+1}) + \theta_2(b_{s+1} - b_s) - (p_{r,s+1} - p_{r+1,s}) \\
  &= -(p_{r+1,s} - p_{r,s}) + (p_{r+1,s+1} - p_{r+1,s}) - p_{r,s+1} + p_{r+1,s} \\
  &= p_{r+1,s+1} - p_{r+1,s} - p_{r,s+1} + p_{r,s} \\
  &= \Delta_F > 0.
\end{align*}
So $y_{10}$ is trapped: locally non-improvable, yet strictly dominated by $y_{01}$.

\smallskip
\emph{Case $\Delta_F < 0$.}
Set $y = y_{00} = (a_r, b_s)$ and choose
\[
  \theta_1 = \frac{p_{r+1,s} - p_{r,s}}{a_{r+1} - a_r} = R_1(y),
  \qquad
  \theta_2 = \frac{p_{r,s+1} - p_{r,s}}{b_{s+1} - b_s} = R_2(y).
\]
Both are nonneg\-ative and lie at the right endpoints of their respective
slope intervals, so $\theta_1 \in [L_1(y), R_1(y)]$ and
$\theta_2 \in [L_2(y), R_2(y)]$.
By Lemma~\ref{lem:slopepin}(b), every adjacent move from $y_{00}$ is
non-improving.  But
\begin{align*}
  u_\theta(y_{11}) - u_\theta(y_{00})
  &= \theta_1(a_{r+1} - a_r) + \theta_2(b_{s+1} - b_s) - (p_{r+1,s+1} - p_{r,s}) \\
  &= (p_{r+1,s} - p_{r,s}) + (p_{r,s+1} - p_{r,s}) - p_{r+1,s+1} + p_{r,s} \\
  &= -(p_{r+1,s+1} - p_{r+1,s} - p_{r,s+1} + p_{r,s}) \\
  &= -\Delta_F > 0.
\end{align*}
So $y_{00}$ is a non-global point with no improving neighbor.

\smallskip
In both cases, the trapped point violates simplicity.
\end{proof}

\begin{proof}[Proof of Theorem~\ref{thm:k2}]
\emph{(Additive $\Rightarrow$ Simple.)}
If $p_0(y_1,y_2) = p_1(y_1) + p_2(y_2)$, we may normalize the summands
as coordinate-line restrictions: $p_1(a_r) = p_0(a_r, b_1) - p_0(a_1, b_1)$
and $p_2(b_s) = p_0(a_1, b_s)$.  Since $p_0$ is convex and nondecreasing,
each $p_i$ is convex and nondecreasing on $M_i$.
The utility then separates, and each coordinate objective is discretely
concave by Lemma~\ref{lem:concavity}.
Therefore every local maximum is global, and any non-maximizer has an
improving adjacent neighbor (this is the $k=2$ case of Lemma~\ref{lem:easy}
below).

\emph{(Simple $\Rightarrow$ Additive.)}
Assume $\Simplenat(Y, p_0)$.
By the contrapositive of Lemma~\ref{lem:trap}, every elementary face
satisfies $\Delta_F = 0$.
Since all adjacent $2 \times 2$ mixed differences vanish, the coordinate-1
increments $p_0(a_{r+1}, b_s) - p_0(a_r, b_s)$ are independent of $b_s$
(moving $b_s$ to an adjacent $b_{s+1}$ changes both terms identically).
By chaining, each increment depends only on $r$.
Fix a basepoint $(a_1, b_1)$ and define
\[
  q_1(a_r) \coloneqq p_0(a_r, b_1) - p_0(a_1, b_1),
  \qquad
  q_2(b_s) \coloneqq p_0(a_1, b_s) - p_0(a_1, b_1).
\]
Telescoping along first coordinate~1 and then coordinate~2 gives
\[
  p_0(a_r, b_s) = p_0(a_1, b_1) + q_1(a_r) + q_2(b_s).
\]
Absorbing $p_0(a_1, b_1)$ into $q_1$ yields the additive decomposition.
\end{proof}

%% ===================================================================
\section{Proof of Theorem~\ref{thm:main} (Main Theorem)}\label{sec:proof}
%% ===================================================================

Although the theorem is stated for $k \geq 2$, we note for completeness
that the case $k=1$ is trivial: additivity is vacuous, and simplicity
follows from discrete concavity (Lemma~\ref{lem:concavity} below).
We assume $k \geq 2$ throughout this section.

%% -------------------------------------------------------------------
\subsection{Easy direction: Additive implies Simple}\label{sec:easy}
%% -------------------------------------------------------------------

\begin{lemma}[Discrete concavity]\label{lem:concavity}
Let $M = \{m_1 < \cdots < m_n\} \subseteq [0,1]$ and let $p\colon M \to \R$
be \emph{convex on the chain} $M$, meaning the secant slopes
\[
  s_r \coloneqq \frac{p(m_{r+1}) - p(m_r)}{m_{r+1} - m_r}, \qquad r = 1,\dots,n-1,
\]
are nondecreasing: $s_1 \leq s_2 \leq \cdots \leq s_{n-1}$.
For $\theta \geq 0$, define $f(m_r) = \theta\, m_r - p(m_r)$.
Then $f$ is \emph{discretely concave} on $M$: every local maximum of~$f$
(with respect to immediate adjacency in $M$) is a global maximum.
\end{lemma}

\begin{proof}
The secant slopes of $f$ are
\[
  \sigma_r \coloneqq \frac{f(m_{r+1}) - f(m_r)}{m_{r+1} - m_r} = \theta - s_r,
  \qquad r = 1,\dots,n-1.
\]
Since the $s_r$ are nondecreasing, the $\sigma_r$ are \emph{nonincreasing}:
$\sigma_1 \geq \sigma_2 \geq \cdots \geq \sigma_{n-1}$.

Suppose $m_{r^*}$ is a local maximum of $f$, i.e.,
$f(m_{r^*}) \geq f(m_{r^*-1})$ (if $r^* > 1$) and
$f(m_{r^*}) \geq f(m_{r^*+1})$ (if $r^* < n$).

\emph{Case 1: $1 < r^* < n$.}
Then $\sigma_{r^*-1} \geq 0$ (from $f(m_{r^*}) \geq f(m_{r^*-1})$, noting $m_{r^*} > m_{r^*-1}$)
and $\sigma_{r^*} \leq 0$ (from $f(m_{r^*}) \geq f(m_{r^*+1})$).
Since the $\sigma_r$ are nonincreasing, we have
$\sigma_r \geq \sigma_{r^*-1} \geq 0$ for all $r < r^*$
and $\sigma_r \leq \sigma_{r^*} \leq 0$ for all $r \geq r^*$.
Therefore, for any $r < r^*$:
\[
  f(m_{r^*}) - f(m_r) = \sum_{t=r}^{r^*-1} \sigma_t (m_{t+1} - m_t) \geq 0,
\]
and for any $r > r^*$:
\[
  f(m_r) - f(m_{r^*}) = \sum_{t=r^*}^{r-1} \sigma_t (m_{t+1} - m_t) \leq 0.
\]
So $f(m_{r^*}) \geq f(m_r)$ for all $r$, i.e., $m_{r^*}$ is a global maximum.

\emph{Case 2: $r^* = 1$.}
Then $\sigma_1 \leq 0$, so $\sigma_r \leq 0$ for all $r \geq 1$,
and $f$ is nonincreasing, making $m_1$ a global maximum.

\emph{Case 3: $r^* = n$.}
Then $\sigma_{n-1} \geq 0$, so $\sigma_r \geq 0$ for all $r \leq n-1$,
and $f$ is nondecreasing, making $m_n$ a global maximum.
\end{proof}

\begin{lemma}[Additive $\Longrightarrow$ Simple]\label{lem:easy}
Suppose $p_0(y) = \sum_{i=1}^k p_i(y_i)$ for all $y \in Y$,
where each $p_i\colon M_i \to \R$ is convex (nondecreasing secant slopes)
and nondecreasing.
Then $\Simplenat(Y, p_0)$.
\end{lemma}

\begin{proof}
For any $\theta \in \Rp^k$, the utility separates:
\[
  u_\theta(y) = \sum_{i=1}^k \bigl[\theta_i y_i - p_i(y_i)\bigr]
  \eqqcolon \sum_{i=1}^k f_i(y_i).
\]
Each $f_i$ is discretely concave on $M_i$ by Lemma~\ref{lem:concavity}
(applied with $p = p_i$ and $\theta = \theta_i$).

Suppose $y \in Y$ is not a global maximizer of $u_\theta$ over $Y$.
Then there exists $y^* \in Y$ with $u_\theta(y^*) > u_\theta(y)$.
Since $u_\theta(y^*) - u_\theta(y) = \sum_{i=1}^k [f_i(y^*_i) - f_i(y_i)] > 0$,
there exists some coordinate $i \in [k]$ with $f_i(y^*_i) > f_i(y_i)$.
In particular, $y_i$ is not a global maximizer of $f_i$ on $M_i$.

By Lemma~\ref{lem:concavity}, $y_i$ is also not a \emph{local} maximum of $f_i$
on $M_i$ (since every local maximum is global).
Therefore, there exists an immediate neighbor $y'_i \in M_i$ of $y_i$
(in the natural order) with $f_i(y'_i) > f_i(y_i)$.

Define $z \in Y$ by $z_i = y'_i$ and $z_j = y_j$ for $j \neq i$.
Then $z \in N_{\mathrm{nat}}(y)$ and
\[
  u_\theta(z) - u_\theta(y) = f_i(y'_i) - f_i(y_i) > 0. \qedhere
\]
\end{proof}

%% -------------------------------------------------------------------
\subsection{Hard direction: Simple implies Additive}\label{sec:hard}
%% -------------------------------------------------------------------

\begin{remark}[False uniform pinning]\label{rem:false}
A natural approach to the slice reduction would be to find, for each frozen
coordinate $\ell \notin \{i,j\}$, a \emph{single} value $\hat\theta_\ell \geq 0$
that pins coordinate $\ell$ at the anchor $a_\ell$ \emph{simultaneously for
all} $y \in S_{ij}(a)$.
That is, one would want
$\hat\theta_\ell \in \bigcap_{y \in S_{ij}(a)} [\sigma_\ell^-(y), \sigma_\ell^+(y)]$,
where $\sigma_\ell^-(y)$ and $\sigma_\ell^+(y)$ are the left and right
discrete slopes of the one-dimensional section $t \mapsto p_0(y|_{\ell \to t})$ at $a_\ell$.

This \textbf{uniform} version is \textbf{false}.
Consider $k = 3$ with
\[
  p_0(y_1, y_2, y_3) = \tfrac{1}{2}(y_1 + y_3)^2,
\]
$M_1 = M_2 = \{0,1\}$, $M_3 = \{0, \tfrac{1}{2}, 1\}$, and anchor $a_3 = \tfrac{1}{2}$.
The one-dimensional section in coordinate~$3$ at anchor $a_3 = \frac{1}{2}$ is
$g(t) = \frac{1}{2}(y_1 + t)^2$.
At $y_1 = 0$: the supporting-slope interval at $t = \frac{1}{2}$ is
$[\sigma^-, \sigma^+] = [\tfrac{1}{4}, \tfrac{3}{4}]$.
At $y_1 = 1$: the supporting-slope interval at $t = \frac{1}{2}$ is
$[\sigma^-, \sigma^+] = [\tfrac{5}{4}, \tfrac{7}{4}]$.
These intervals are \emph{disjoint}, so no single $\hat\theta_3$ pins
coordinate~$3$ at $a_3 = \frac{1}{2}$ for all $y \in S_{12}(a)$ simultaneously.

\smallskip\noindent
\emph{Why joint convexity does not rescue this.}
One might hope that the monotonicity of the subdifferential
$\partial p_0$ (which holds by joint convexity) would force the slope intervals to overlap.
However, the relevant difference vectors have zero $\ell$-component:
when $y$ and $y'$ differ only in coordinates $i,j$ (i.e., within the slice $S_{ij}(a)$),
the vector $y - y'$ has $(y-y')_\ell = 0$, so subdifferential monotonicity gives
$\langle \xi - \xi', y - y' \rangle \geq 0$
with the $\ell$-components of $\xi, \xi'$ multiplied by zero.
This provides no constraint on the $\ell$-th slope intervals.
\end{remark}

\begin{lemma}[Global pointwise pinning]\label{lem:pointwise}
Fix $i \neq j$, an anchor $a \in \prod_{\ell \neq i,j} M_\ell$,
a point $y \in S_{ij}(a)$, and a frozen coordinate $\ell \notin \{i,j\}$.
There exists $\hat\theta_\ell(y) \geq 0$ such that for every
$t \in M_\ell$,
\[
  \hat\theta_\ell(y)\, t - p_0(y|_{\ell \to t})
  \;\leq\;
  \hat\theta_\ell(y)\, a_\ell - p_0(y).
\]
That is, no move in coordinate $\ell$---adjacent or not---can strictly
improve the utility component in that coordinate.
\end{lemma}

\begin{proof}
Consider the one-dimensional section
\[
  g(t) \coloneqq p_0(y|_{\ell \to t}), \qquad t \in M_\ell.
\]
Since $p_0$ is convex on $[0,1]^k$, the restriction $g$ is convex on the
chain $M_\ell$ (secant slopes nondecreasing).
Since $p_0$ is coordinatewise nondecreasing, $g$ is nondecreasing.

We need $\hat\theta_\ell \geq 0$ such that $\phi(t) \coloneqq \hat\theta_\ell\, t - g(t)$
is maximized at $t = a_\ell$ over \emph{all} of $M_\ell$.

Let $a_\ell = m_{\ell,q}$, i.e., $a_\ell$ is the $q$-th element of $M_\ell$.
Define the left and right secant slopes (when they exist):
\[
  \sigma^- \coloneqq \frac{g(m_{\ell,q}) - g(m_{\ell,q-1})}{m_{\ell,q} - m_{\ell,q-1}}, \qquad
  \sigma^+ \coloneqq \frac{g(m_{\ell,q+1}) - g(m_{\ell,q})}{m_{\ell,q+1} - m_{\ell,q}}.
\]
By convexity of $g$, $\sigma^- \leq \sigma^+$ whenever both are defined.
By monotonicity of $g$, both slopes are nonnegative whenever defined.

\emph{Case 1: Interior point ($1 < q < n_\ell$).}
Choose $\hat\theta_\ell \in [\sigma^-, \sigma^+]$ (the interval is nonempty by convexity).
For the left neighbor $m_{\ell,q-1}$:
\[
  \hat\theta_\ell (m_{\ell,q} - m_{\ell,q-1}) - [g(m_{\ell,q}) - g(m_{\ell,q-1})]
  = (\hat\theta_\ell - \sigma^-)(m_{\ell,q} - m_{\ell,q-1}) \geq 0,
\]
so the move from $a_\ell$ to $m_{\ell,q-1}$ is weakly non-improving.
For the right neighbor $m_{\ell,q+1}$:
\[
  \hat\theta_\ell (m_{\ell,q+1} - m_{\ell,q}) - [g(m_{\ell,q+1}) - g(m_{\ell,q})]
  = (\hat\theta_\ell - \sigma^+)(m_{\ell,q+1} - m_{\ell,q}) \leq 0,
\]
so the move from $a_\ell$ to $m_{\ell,q+1}$ is also weakly non-improving.

\emph{Case 2: Left boundary ($q = 1$).}
The only adjacent move is rightward, to $m_{\ell,2}$.
Set $\hat\theta_\ell = 0$.
Then the utility change from moving to $m_{\ell,2}$ is
\[
  0 \cdot m_{\ell,2} - g(m_{\ell,2}) - [0 \cdot m_{\ell,1} - g(m_{\ell,1})]
  = g(m_{\ell,1}) - g(m_{\ell,2}) \leq 0,
\]
since $g$ is nondecreasing.

\emph{Case 3: Right boundary ($q = n_\ell$).}
The only adjacent move is leftward, to $m_{\ell,n_\ell-1}$.
Set $\hat\theta_\ell = \sigma^-$, the left secant slope (which is $\geq 0$).
The utility change from moving to $m_{\ell,n_\ell-1}$ is
\[
  \sigma^-(m_{\ell,n_\ell-1} - m_{\ell,n_\ell}) - [g(m_{\ell,n_\ell-1}) - g(m_{\ell,n_\ell})]
  = (\sigma^- - \sigma^-)(m_{\ell,n_\ell-1} - m_{\ell,n_\ell}) = 0 \leq 0,
\]
so the move is weakly non-improving.
In all cases, $\hat\theta_\ell \geq 0$ and $a_\ell$ is a local maximizer of
$\phi(t) = \hat\theta_\ell\, t - g(t)$ among the natural-adjacent values.
Since $g$ is convex, $\phi$ is discretely concave on $M_\ell$.
By Lemma~\ref{lem:concavity}, every local maximizer of a discretely
concave function is a global maximizer.
Therefore $\phi(a_\ell) \geq \phi(t)$ for \emph{all} $t \in M_\ell$,
not merely for adjacent~$t$.
\end{proof}

\begin{lemma}[Slice reduction]\label{lem:slice}
Let $k \geq 2$ and assume $\Simplenat(Y, p_0)$.
Then for all distinct $i, j \in [k]$ and every anchor $a \in \prod_{\ell \neq i,j} M_\ell$,
the two-coordinate slice $S_{ij}(a)$ equipped with $p_0|_{S_{ij}(a)}$ and the
natural orders on $M_i, M_j$ is simple.
\end{lemma}

\begin{proof}
Fix $(\theta_i, \theta_j) \in \Rp^2$ and $y \in S_{ij}(a)$ such that $y$ is not
a global maximizer of
\[
  v_{(\theta_i,\theta_j)}(y) \coloneqq \theta_i y_i + \theta_j y_j - p_0(y)
\]
over $S_{ij}(a)$.  We must produce $z \in N_{\mathrm{nat}}(y) \cap S_{ij}(a)$
with $v_{(\theta_i,\theta_j)}(z) > v_{(\theta_i,\theta_j)}(y)$.

\emph{Step 1: Construct the extended type.}
For each frozen coordinate $\ell \notin \{i,j\}$,
apply Lemma~\ref{lem:pointwise} with the fixed point $y$ to obtain
$\hat\theta_\ell(y) \geq 0$ that pins coordinate $\ell$ at $a_\ell$.
Define the full type $\theta = (\theta_i, \theta_j, \hat\theta_{-\{i,j\}}(y)) \in \Rp^k$.
(The pinning coordinates depend on $y$; no uniform choice over the slice
is claimed or needed --- cf.\ Remark~\ref{rem:false}.)

\emph{Step 2: Non-optimality on the full menu.}
There exists $w \in S_{ij}(a)$ with $v_{(\theta_i,\theta_j)}(w) > v_{(\theta_i,\theta_j)}(y)$.
Since $w$ and $y$ agree on all frozen coordinates, we have
\[
  u_\theta(w) - u_\theta(y) = v_{(\theta_i,\theta_j)}(w) - v_{(\theta_i,\theta_j)}(y)
  + \sum_{\ell \notin \{i,j\}} \hat\theta_\ell(y)(w_\ell - y_\ell) = v_{(\theta_i,\theta_j)}(w) - v_{(\theta_i,\theta_j)}(y) > 0,
\]
where the sum vanishes because $w_\ell = a_\ell = y_\ell$ for all frozen $\ell$.
So $y$ is not a global maximizer of $u_\theta$ on $Y$.

\emph{Step 3: Apply full-menu simplicity.}
By $\Simplenat(Y, p_0)$, there exists $z \in N_{\mathrm{nat}}(y)$ with
$u_\theta(z) > u_\theta(y)$.
The neighbor $z$ differs from $y$ in exactly one coordinate.

\emph{Step 4: The neighbor must lie in the slice.}
Suppose, for contradiction, that $z$ differs from $y$ in some frozen
coordinate $\ell \notin \{i,j\}$, so $z_\ell \neq y_\ell = a_\ell$
and $z_{-\ell} = y_{-\ell}$.
Then
\[
  u_\theta(z) - u_\theta(y)
  = \hat\theta_\ell(y)(z_\ell - a_\ell) - [p_0(y|_{\ell \to z_\ell}) - p_0(y)] \leq 0,
\]
by the pinning property of $\hat\theta_\ell(y)$ (Lemma~\ref{lem:pointwise}).
This contradicts $u_\theta(z) > u_\theta(y)$.
Therefore $z \in S_{ij}(a)$, and $z$ differs from $y$ in coordinate $i$ or $j$ only.
Since $u_\theta(z) - u_\theta(y) = v_{(\theta_i,\theta_j)}(z) - v_{(\theta_i,\theta_j)}(y) > 0$,
the point $z$ is the required improving neighbor within the slice.
\end{proof}

\begin{lemma}[Pairwise vanishing implies additivity]\label{lem:vanishing}
Let $k \geq 2$.  Suppose that for every pair of distinct coordinates $i, j \in [k]$
and every anchor $a \in \prod_{\ell \neq i,j} M_\ell$,
all adjacent $2 \times 2$ minors of the slice price matrix vanish:
\[
  P^{ij,a}_{r+1,s+1} - P^{ij,a}_{r+1,s} - P^{ij,a}_{r,s+1} + P^{ij,a}_{r,s} = 0
\]
for all $r \in [n_i - 1]$, $s \in [n_j - 1]$.
Then $p_0$ is additively separable on $Y$: there exist functions $q_\ell\colon M_\ell \to \R$
such that
\[
  p_0(y) = \sum_{\ell=1}^k q_\ell(y_\ell) \qquad \forall\, y \in Y.
\]
\end{lemma}

\begin{proof}
\emph{Step 1: Coordinate increments are independent of other coordinates.}
For each coordinate $i \in [k]$ and index $r \in [n_i - 1]$, define the \emph{increment}
\[
  \Delta_{i,r}(x_{-i}) \coloneqq p_0(x|_{i \to m_{i,r+1}}) - p_0(x|_{i \to m_{i,r}}),
\]
where $x_{-i}$ denotes the values of all coordinates except~$i$.
We claim that $\Delta_{i,r}(x_{-i})$ depends only on $r$, not on $x_{-i}$.

Fix any $j \neq i$.  The vanishing $2\times 2$ minor condition for the pair $(i,j)$
at anchor $a = x_{-\{i,j\}}$ states:
\begin{align*}
  &p_0(x|_{i \to m_{i,r+1},\, j \to m_{j,s+1}}) - p_0(x|_{i \to m_{i,r},\, j \to m_{j,s+1}}) \\
  &\qquad = p_0(x|_{i \to m_{i,r+1},\, j \to m_{j,s}}) - p_0(x|_{i \to m_{i,r},\, j \to m_{j,s}})
\end{align*}
for all $s \in [n_j - 1]$.
In other words, $\Delta_{i,r}$ is unchanged when $x_j$ moves from $m_{j,s}$ to $m_{j,s+1}$.
By chaining over adjacent steps in $M_j$, $\Delta_{i,r}$ is constant across all values of~$x_j$.
Since this holds for every $j \neq i$, the increment $\Delta_{i,r}$ is independent of
$x_{-i}$ entirely.  Write $d_{i,r} \coloneqq \Delta_{i,r}$.

\emph{Step 2: Basepoint telescoping.}
Fix a basepoint $y^* \in Y$.  For each $\ell \in [k]$, define
\[
  q_\ell(y_\ell) \coloneqq p_0(y^*|_{\ell \to y_\ell}) - p_0(y^*).
\]
Note $q_\ell(y^*_\ell) = 0$ for all $\ell$.
For an arbitrary $y \in Y$, telescope along the path that changes one coordinate at a time:
\[
  y^* \;\to\; (y_1, y^*_2, \dots, y^*_k) \;\to\; (y_1, y_2, y^*_3, \dots, y^*_k) \;\to\; \cdots \;\to\; y.
\]
The total change is
\[
  p_0(y) - p_0(y^*) = \sum_{\ell=1}^k \bigl[p_0(y_1,\dots,y_\ell, y^*_{\ell+1}, \dots, y^*_k) -
  p_0(y_1,\dots,y_{\ell-1}, y^*_\ell, y^*_{\ell+1},\dots,y^*_k)\bigr].
\]
The $\ell$-th summand is the cost of changing coordinate $\ell$ from $y^*_\ell$ to $y_\ell$.
By Step~1, this cost depends only on $y^*_\ell$, $y_\ell$, and the increments $d_{\ell,r}$,
and is independent of the current values of the other coordinates.
Concretely, if $y^*_\ell = m_{\ell,a}$ and $y_\ell = m_{\ell,b}$ with $b > a$, then
\[
  p_0(\cdots|_{\ell \to m_{\ell,b}}) - p_0(\cdots|_{\ell \to m_{\ell,a}}) = \sum_{r=a}^{b-1} d_{\ell,r},
\]
which depends only on $y^*_\ell$ and $y_\ell$ (and the case $b < a$ is analogous).
In particular, the $\ell$-th summand equals $q_\ell(y_\ell) - q_\ell(y^*_\ell) = q_\ell(y_\ell)$.
Therefore
\[
  p_0(y) = p_0(y^*) + \sum_{\ell=1}^k q_\ell(y_\ell).
\]
Absorbing the constant $p_0(y^*)$ into $q_1$ completes the proof.
\end{proof}

\begin{lemma}[Assembly]\label{lem:assembly}
If $\Simplenat(Y, p_0)$, then $p_0$ is additively separable on $Y$.
\end{lemma}

\begin{proof}
We chain the preceding results:
\begin{enumerate}[label=(\arabic*)]
  \item By Lemma~\ref{lem:slice} (slice reduction), for every pair of distinct coordinates
    $i,j \in [k]$ and every anchor $a \in \prod_{\ell \neq i,j} M_\ell$,
    the slice $S_{ij}(a)$ equipped with $p_0|_{S_{ij}(a)}$ and natural orders is simple.
  \item By Theorem~\ref{thm:k2} (the $k=2$ characterization), simplicity
    of $S_{ij}(a)$ implies the slice price matrix $P^{ij,a}$ is additively separable.
    Additive separability of a two-index matrix is equivalent to the vanishing of all
    adjacent $2 \times 2$ minors:
    \[
      P^{ij,a}_{r+1,s+1} - P^{ij,a}_{r+1,s} - P^{ij,a}_{r,s+1} + P^{ij,a}_{r,s} = 0
      \quad \forall\, r \in [n_i - 1],\; s \in [n_j - 1].
    \]
  \item Since this holds for every pair $(i,j)$ and every anchor $a$,
    the hypotheses of Lemma~\ref{lem:vanishing} are satisfied.
  \item By Lemma~\ref{lem:vanishing}, $p_0$ is additively separable on~$Y$.  \qedhere
\end{enumerate}
\end{proof}

\begin{proof}[Proof of Theorem~\ref{thm:main}]
The forward implication ($\Simplenat \Rightarrow$ additive) is Lemma~\ref{lem:assembly}.
For the reverse, suppose $p_0(y) = \sum_i p_i(y_i)$ on $Y$.
Since all adjacent $2 \times 2$ minors then vanish, we may normalize the
summands via basepoint telescoping as in Lemma~\ref{lem:vanishing}:
$q_\ell(y_\ell) = p_0(y^*|_{\ell \to y_\ell}) - p_0(y^*)$.
Each $q_\ell$ is a restriction of $p_0$ to a coordinate line through $y^*$,
hence convex and nondecreasing on $M_\ell$.
The hypotheses of Lemma~\ref{lem:easy} are therefore satisfied,
and $\Simplenat(Y, p_0)$ follows.
\end{proof}

%% ===================================================================
\section{Proof of Theorem~\ref{thm:order-norm} (Order Normalization)}\label{sec:order-norm}
%% ===================================================================

\begin{lemma}[Order normalization]\label{lem:order-norm}
If there exists a family of total orders $\preceq = (\preceq_i)_{i=1}^k$
on the sets $M_i$ such that $\Simpleord(Y, p_0)$, then $\Simplenat(Y, p_0)$.
\end{lemma}

\begin{proof}
We proceed in four steps.

\emph{Step 1: Slice reduction under $\preceq$.}
Fix distinct $i,j \in [k]$ and an anchor $a \in \prod_{\ell \neq i,j} M_\ell$.
The argument of Lemma~\ref{lem:slice} goes through verbatim with
$\preceq$-neighbors in place of nat-neighbors:
for each $y \in S_{ij}(a)$ and each frozen $\ell \notin \{i,j\}$,
Lemma~\ref{lem:pointwise} provides $\hat\theta_\ell(y) \geq 0$ pinning
coordinate $\ell$ at $a_\ell$.
By Lemma~\ref{lem:pointwise}, the pinning inequality
$\hat\theta_\ell(y)\, t - p_0(y|_{\ell \to t}) \leq
 \hat\theta_\ell(y)\, a_\ell - p_0(y)$
holds for \emph{every} $t \in M_\ell$, not just adjacent values.
In particular, any $\preceq_\ell$-adjacent move in a frozen coordinate
is non-improving.
Therefore, if $y$ is not optimal on the slice and the full type is
$\theta = (\theta_i, \theta_j, \hat\theta_{-\{i,j\}}(y))$,
then $\Simpleord(Y,p_0)$ produces an improving $\preceq$-neighbor
that must lie in $S_{ij}(a)$ (frozen coordinates are pinned).
Thus every slice $S_{ij}(a)$ is simple under $(\preceq_i, \preceq_j)$.

\emph{Step 2: Contrapositive --- non-additive slice yields a contradiction.}
Suppose for contradiction that some slice $S_{ij}(a)$ has a non-additive price matrix.
By the contrapositive of Theorem~\ref{thm:k2},
there exist $\eta = (\eta_i, \eta_j) \in \Rp^2$ and $y \in S_{ij}(a)$
such that $y$ is not a global maximizer of
$v_\eta(w) = \eta_i w_i + \eta_j w_j - p_0(w)$ on $S_{ij}(a)$,
yet $y$ has \emph{no} nat-improving neighbor in $S_{ij}(a)$.
This means $y_i$ is a local maximum of $w_i \mapsto \eta_i w_i - p_0(w|_{i \to w_i})$
on $M_i$ (holding $y_j$ and $a$ fixed), and similarly $y_j$ is a local maximum
of the corresponding function on $M_j$.

By Lemma~\ref{lem:concavity} (discrete concavity), both local maxima are
\emph{global} maxima on their respective coordinate lines.
Thus the absence of a nat-improving neighbor forces $y_i$ and $y_j$ to be
global maximizers of their coordinate-line utilities, so no
one-coordinate move---regardless of which order defines adjacency---can
improve $v_\eta$.

Now consider improvement under any adjacency order $\preceq$.
Any $\preceq$-neighbor $z \in N_\preceq(y) \cap S_{ij}(a)$ differs from $y$
in exactly one of $\{i,j\}$.  If $z$ differs in coordinate~$i$, then
$v_\eta(z) - v_\eta(y) = [\eta_i z_i - p_0(y|_{i \to z_i})] - [\eta_i y_i - p_0(y)] \leq 0$
since $y_i$ is a global maximizer.  Similarly if $z$ differs in coordinate~$j$.
So $y$ has no $\preceq$-improving neighbor in $S_{ij}(a)$ either,
contradicting the simplicity of $S_{ij}(a)$ under $(\preceq_i, \preceq_j)$
established in Step~1.

\emph{Step 3: All slice price matrices are additive.}
Step~2 shows that every slice $S_{ij}(a)$ has an additive price matrix,
so all adjacent $2 \times 2$ minors vanish for every pair $(i,j)$ and every anchor $a$.

\emph{Step 4: Conclude.}
By Lemma~\ref{lem:vanishing}, $p_0$ is additively separable on $Y$.
By Lemma~\ref{lem:easy}, $\Simplenat(Y, p_0)$ holds.
\end{proof}

\begin{proof}[Proof of Theorem~\ref{thm:order-norm}]
The chain of implications is:
\[
  \bigl(\exists \preceq\;\; \Simpleord(Y,p_0)\bigr)
  \;\;\xRightarrow{\text{Lem.~\ref{lem:order-norm}}}\;\;
  \Simplenat(Y,p_0)
  \;\;\xRightarrow{\text{Lem.~\ref{lem:assembly}}}\;\;
  p_0 \text{ additive on } Y
  \;\;\xRightarrow{\text{Lem.~\ref{lem:easy}}}\;\;
  \Simplenat(Y,p_0).
\]
The converse $\Simplenat \Rightarrow (\exists \preceq\;\;\Simpleord)$ is trivial,
since $\Simplenat$ is itself simplicity under the natural orders.
\end{proof}

\begin{remark}[Order normalization under the two neighbor notions]%
\label{rem:order-norm-scope}
Theorem~\ref{thm:order-norm} has genuine content because we use
\emph{immediate-adjacency} neighbors: different total orders on a
coordinate set $M_i$ produce different adjacency graphs
(e.g., on $\{0, 0.5, 1\}$, the natural order makes $0$ and $1$
non-adjacent, while the order $0 \prec 1 \prec 0.5$ makes them
adjacent).  The theorem says that among all such adjacency structures,
the natural numerical ordering already suffices.

Under the original model's neighbor definition
(Definition~\ref{def:broad-simplicity}), order normalization is
\emph{automatic}: since every pair in a total order is comparable,
the neighbor set $N_\preceq(y)$ is the same for every choice of
$\preceq_i$, and coincides with the broad-neighbor set.
Thus the order-selection clause in Theorem~\ref{thm:order-norm} is
vacuous under the original definition --- the orders do no mathematical
work.  This is not a defect; it simply means the original definition
already possesses the order-independence property for free, and
Theorem~\ref{thm:order-norm} establishes the analogous property for
the finer-grained immediate-adjacency formulation.
\end{remark}

%% ===================================================================
\section{Discussion}\label{sec:discussion}
%% ===================================================================

\paragraph{Summary.}
Theorem~\ref{thm:main} establishes that on a finite product menu
$Y = \prod_{i=1}^k M_i$, the simplicity property (every suboptimal bundle admits
an improving immediate neighbor) holds if and only if the canonical price
$p_0$ is additively separable on $Y$.
Theorem~\ref{thm:order-norm} further shows an order-normalization
result under the immediate-adjacency formulation: if simplicity holds
for \emph{some} family of total orders on the coordinate sets, then it
also holds under the natural numerical order.
Since different orderings produce genuinely different adjacency graphs,
this says the natural orders are canonical among successful orderings.
(Under the original model's broad-neighbor definition, order normalization
is automatic; see Remark~\ref{rem:order-norm-scope}.)

The proof strategy reduces the $k$-good problem to the
$k=2$ case (Theorem~\ref{thm:k2}) via a sequence of steps:
\begin{enumerate}[label=(\roman*)]
  \item \emph{Two-good trap construction} (Lemma~\ref{lem:trap}): a nonzero
    $2 \times 2$ mixed difference creates a point that is locally non-improvable
    yet globally suboptimal, by choosing secant slopes that simultaneously pin
    both coordinates.
  \item \emph{Pointwise pinning} (Lemma~\ref{lem:pointwise}): for each
    fixed point in a two-coordinate slice, frozen coordinates can be
    pinned using one-dimensional supporting slopes from convexity.
  \item \emph{Slice reduction} (Lemma~\ref{lem:slice}): simplicity of the
    full menu induces simplicity of every two-coordinate slice.
  \item \emph{$k=2$ theorem} (Theorem~\ref{thm:k2}): each slice has additive price.
  \item \emph{Discrete integration} (Lemma~\ref{lem:vanishing}): pairwise
    additive separability implies global additive separability.
\end{enumerate}

\paragraph{The false uniform-pinning lemma as a cautionary note.}
As documented in Remark~\ref{rem:false}, the \emph{uniform} pinning statement
--- choosing a single supporting slope $\hat\theta_\ell$ that works
simultaneously for all points in a slice --- is false.
The explicit counterexample ($p_0 = \frac{1}{2}(y_1+y_3)^2$ with
$M_3 = \{0,\frac{1}{2},1\}$) shows that joint convexity of $p_0$ does not
force the supporting-slope intervals at a given anchor to intersect across
all slice points.
The correct replacement is the \emph{pointwise} version
(Lemma~\ref{lem:pointwise}), which suffices because the slice reduction
fixes the suboptimal point $y$ \emph{before} choosing the extended type.
The quantifier order --- $\forall y$, $\exists \hat\theta(y)$ --- is
the crucial distinction.

\paragraph{Open question: beyond convex $p_0$.}
The entire argument relies on the convexity of $p_0$ (inherited from the
IC/IR/NPT mechanism-design model) in three places:
discrete concavity (Lemma~\ref{lem:concavity}),
pointwise pinning (Lemma~\ref{lem:pointwise}), and
the monotonicity of secant slopes used in the $k=2$ theorem.
It is an open question whether the equivalence between simplicity and
additive separability extends to non-convex pricing functions on product grids.

\paragraph{Out of scope.}
This draft addresses only finite product menus with convex canonical
pricing.  The extension to infinite (or continuous) menus is deferred.
As shown in Lemma~\ref{lem:broad-narrow} and
Remark~\ref{rem:model-pdf}, the original model's neighbor definition
(any one-coordinate change) coincides with immediate-adjacency
simplicity in the finite convex setting, so our characterization
covers the original model's notion on its intended domain.
What remains open is whether the equivalence extends beyond the
finite-grid convex regime.

\end{document}
